pdfoutput=1
\pdfoutput=1
\documentclass[12pt,a4paper]{article}

% ============================================================
%  CD-Theorem — Causal Discovery Theorem:
%  On the Separability of Label Noise and Unobserved Confounding
%  arXiv-ready · English · Standalone (SCX-citable)
% ============================================================

\usepackage[T1]{fontenc}
\usepackage{amsmath,amssymb,amsfonts}
\usepackage{mathtools}
\usepackage{amsthm}
\usepackage{bm}
\usepackage{graphicx}
\usepackage{xcolor}
\usepackage{hyperref}
\usepackage{booktabs}
\usepackage{enumitem}
\usepackage[numbers]{natbib}
\usepackage{dsfont}
\usepackage{bbm}
\usepackage{geometry}
\geometry{left=2.5cm,right=2.5cm,top=2.5cm,bottom=2.5cm}

% ---------- Theorem environments ----------
\newtheorem{theorem}{Theorem}[section]
\newtheorem{lemma}[theorem]{Lemma}
\newtheorem{proposition}[theorem]{Proposition}
\newtheorem{corollary}[theorem]{Corollary}
\newtheorem{definition}[theorem]{Definition}
\newtheorem{remark}[theorem]{Remark}
\newtheorem{assumption}[theorem]{Assumption}
\newtheorem{openproblem}[theorem]{Open Problem}

% ---------- Status tags ----------
\newcommand{\rigorous}{\textcolor{green!50!black}{\small\textsf{[Rigorous]}}}
\newcommand{\conditionallyrigorous}{\textcolor{orange}{\small\textsf{[Conditionally Rigorous]}}}
\newcommand{\openquest}{\textcolor{red!70!black}{\small\textsf{[Open]}}}

% ---------- Honest-critique environment ----------
\newenvironment{critique}{%
    \medskip\noindent
    {\bfseries\sffamily\color{red!60!black} Honest Critique --- }
    \itshape\small
}{\par\medskip}

% ---------- Macros ----------
\newcommand{\R}{\mathbb{R}}
\newcommand{\N}{\mathbb{N}}
\newcommand{\E}{\mathbb{E}}
\newcommand{\V}{\mathbb{V}}
\newcommand{\I}{\mathbb{I}}
\newcommand{\cX}{\mathcal{X}}
\newcommand{\cY}{\mathcal{Y}}
\newcommand{\cD}{\mathcal{D}}
\newcommand{\cF}{\mathcal{F}}
\newcommand{\cG}{\mathcal{G}}
\newcommand{\cA}{\mathcal{A}}
\newcommand{\cI}{\mathcal{I}}
\newcommand{\cM}{\mathcal{M}}
\newcommand{\cH}{\mathcal{H}}
\newcommand{\cP}{\mathcal{P}}
\newcommand{\cB}{\mathcal{B}}
\newcommand{\ind}[1]{\mathbf{1}_{[#1]}}
\newcommand{\argmax}{\operatorname*{arg\,max}}
\newcommand{\argmin}{\operatorname*{arg\,min}}
\newcommand{\KL}{D_{\mathrm{KL}}}
\newcommand{\TV}{\mathrm{TV}}
\newcommand{\Situs}{\mathrm{Situs}}
\newcommand{\Cercis}{\mathrm{Cercis}}
\newcommand{\supp}{\mathrm{supp}}
\newcommand{\tr}{\mathrm{tr}}
\newcommand{\curl}{\nabla\! \times\!}
\newcommand{\grad}{\nabla}
\newcommand{\Hess}{\mathrm{H}}
\newcommand{\Law}{\mathcal{L}}
\newcommand{\plim}{\xrightarrow{p}}
\newcommand{\dlim}{\xrightarrow{d}}
\newcommand{\indist}{\stackrel{d}{=}}
\newcommand{\defeq}{\triangleq}
\newcommand{\iid}{\mathrm{i.i.d.}}
\newcommand{\bias}{\mathrm{bias}}
\newcommand{\var}{\mathrm{var}}
\newcommand{\mse}{\mathrm{MSE}}
\newcommand{\rank}{\mathrm{rank}}

\title{\bfseries On the Separability of Label Noise and Unobserved\\
Confounding in the SCX Causal Discovery Framework\\
\large --- Strictly Formalized Version ---}

\author{SCX}

\date{\today}

\begin{document}
\maketitle

\begin{abstract}
This document presents the strictly formalized version of the separability theorem for label noise and unobserved confounding within the \textsf{SCX} causal discovery framework. Under the \textsf{SCX} notational conventions, we provide rigorous mathematical statements, enumerated premise assumptions, complete proofs, rigor annotations, and honest critique analyses for three core theorems:
(1)~Strong Causal Inseparability Theorem (Theorem CD.1): when the Sprint parameter $\eta$ is independent of unobserved confounders $z$, label noise and causal confounding are absolutely indistinguishable, with test power bounded by the significance level;
(2)~Conditional Weak Separability Theorem (Theorem CD.2): when $\I(\eta;z)>0$ and the Situs encoding preserves the causal skeleton, a curl-statistic-based test achieves non-trivial power;
(3)~Causal Relay Bound Theorem (Theorem CD.3): under the front-door criterion and T7 cross-domain fidelity, the Cercis bias induced by confounding is bounded by a Situs-distance-dependent term plus statistical error.
\end{abstract}

% ============================================================
\section{Introduction and Formal Setup}
\label{sec:intro}
% ============================================================

\subsection{Problem Setting}

Consider an audit dataset $\cD = \{(x_i, y_i)\}_{i=1}^n$, with feature vectors $x_i \in \cX \subseteq \R^d$ and labels $y_i \in \cY \subseteq \R$. The label generation process is
\begin{equation}\label{eq:label-process}
    y_i = f^*(x_i) + \varepsilon_i + \beta \cdot z_i,
\end{equation}
where $f^*: \cX \to \cY$ is the true labeling function, $\varepsilon_i \sim P_\varepsilon$ is independent label noise (zero mean, finite variance), $z_i \in \R^m$ is an unobserved confounding variable, $\beta \in \R^m$ is the confounding strength vector, and $\beta\cdot z$ denotes the Euclidean inner product.

The auditor's observable information set is
\begin{equation}\label{eq:info-set}
    \cI_A = \{S(x_i),\; \nabla S(x_i),\; H_S(x_i)\}_{i=1}^n,
\end{equation}
where $S(x)$ is the Cercis Score, $\nabla S(x)$ is its gradient, and $H_S(x)$ is its Hessian matrix. The auditor \textbf{cannot} observe the decomposition into $\varepsilon$ and $z$.

\subsection{Core Research Question}

We study the following hypothesis testing problem:
\begin{equation}\label{eq:testing-problem}
    H_0: \beta = 0 \quad (\text{pure label noise})\quad \text{vs.}\quad
    H_1: \beta \neq 0 \quad (\text{confounding present}).
\end{equation}
The goal is to determine: does there exist a test function $\psi: \cI_A \to \{0,1\}$ such that the power $\E_{H_1}[\psi] > \E_{H_0}[\psi] = \alpha$?

\subsection{Contributions}

\begin{enumerate}[label=(\arabic*)]
    \item \textbf{Strong Inseparability} (Theorem~\ref{thm:strong-insep}):
          When $\eta \perp z$, for any level-$\alpha$ test $\psi$,
          $\E_{H_1}[\psi] \leq \alpha + O(n^{-1/2})$.
          We construct two observationally equivalent worlds (World A: confounding, World B: pure noise),
          rigorously proving the likelihood ratio equals 1.
          \rigorous
    \item \textbf{Conditional Weak Separability} (Theorem~\ref{thm:weak-sep}):
          When $\I(\eta;z) > 0$ and Situs preserves the causal skeleton,
          there exists a statistic $\tau$ such that $\E_{H_1}[\tau] > \E_{H_0}[\tau]$,
          achieving weak separation in the Chebyshev inequality sense.
          This result depends conditionally on the Situs skeleton preservation assumption.
          \conditionallyrigorous~/ \openquest~(unconditional form)
    \item \textbf{Causal Relay Bound} (Theorem~\ref{thm:causal-relay}):
          Under the front-door criterion and T7 cross-domain fidelity,
          $\|\Delta S\|_{\cD_t} \leq L_{\Situs}\cdot\varepsilon + \kappa\cdot\delta + C\cdot\varepsilon\cdot\delta + O(1/\sqrt{n})$.
          Each error term is explicitly computable.
          \rigorous
\end{enumerate}

% ============================================================
\section{Preliminaries and Notation}
\label{sec:prelim}
% ============================================================

\subsection{SCX Framework Essentials}

\paragraph{Cercis Score.}
The Cercis Score is a scalar field on $\cX$:
\begin{equation}\label{eq:cercis-def}
    S(x) = Q(x) + \eta N(x),\qquad x \in \cX,
\end{equation}
where $Q(x) \in [0,1]$ is label quality (expert agreement / model confidence),
$N(x) \in [0,1]$ is novelty (dissimilarity to previously audited samples),
and $\eta \geq 0$ is the Sprint noise--difficulty coupling coefficient.

\paragraph{Situs Encoding.}
The Situs operator encodes a data-generating distribution $P(X,Y)$ into a metric-measure space:
\begin{equation}\label{eq:situs-def}
    \Situs(P) = (\cX, d_P, \mu_P),
\end{equation}
where $d_P: \cX \times \cX \to \R_{\geq 0}$ is a task-adapted metric and $\mu_P$ is the marginal measure on $\cX$.
The Situs distance between two distributions is defined as
\begin{equation}\label{eq:situs-dist}
    d_{\Situs}(P, Q) \defeq W_1\big(\Situs(P),\, \Situs(Q)\big),
\end{equation}
the 1-Wasserstein distance between their Situs encodings.

\paragraph{Theorem~3 (Noise--Difficulty Unidentifiability).}
Let $\cA$ be any algorithm that, given $\{S(x_i)\}_{i=1}^n$, outputs a binary classification $\hat{c}_i \in \{\text{noise},\text{hard}\}$ for each sample.
Under \textsf{SCX} assumptions A1--A6, for any significance level $\alpha \in (0,1)$:
\begin{equation}\label{eq:thm3-statement}
    \sup_{\cA}\; \big|\E[\hat{c}_i = \text{noise} \mid \text{true noise}]
    - \E[\hat{c}_i = \text{noise} \mid \text{true hard}]\big|
    \;\leq\; \alpha + o_n(1),
\end{equation}
where $o_n(1) \to 0$ as $n \to \infty$.

\subsection{Causal Graph Notation}

Let $\cG = (V, E)$ be a directed acyclic graph (DAG) with vertex set $V = \{X, Y, Z, U\}$:
\begin{itemize}
    \item $X$: observed features
    \item $Y$: labels
    \item $Z$: unobserved confounders (affecting both $X$ and $Y$)
    \item $U$: all other unobserved variables
\end{itemize}
We assume the causal Markov condition and faithfulness with respect to $\cG$.

\section{Strong Causal Inseparability Theorem (Theorem CD.1)}
\label{sec:strong}

\begin{theorem}[Strong Causal Inseparability]\label{thm:strong-insep}
    Let $\cD = \{(x_i, y_i)\}_{i=1}^n$ be generated according to~\eqref{eq:label-process},
    with $z_i$ unobserved confounders. The auditor's information set is $\cI_A = \{S(x_i), \nabla S(x_i), H_S(x_i)\}_{i=1}^n$.
    If the Sprint noise parameter satisfies $\eta \perp z$ (independence), then for any test function
    $\psi: \cI_A \to \{0,1\}$ testing
    \begin{equation*}
        H_0: \beta = 0\;(\text{pure noise}) \quad\text{vs.}\quad
        H_1: \beta \neq 0\;(\text{confounding present}),
    \end{equation*}
    the power function satisfies
    \begin{equation}\label{eq:power-bound}
        \sup_{\psi}\; \E_{H_1}[\psi] \;\leq\; \alpha + O(n^{-1/2}),
    \end{equation}
    where $\alpha = \E_{H_0}[\psi]$ is the test level.
    That is: no test can distinguish label noise from unobserved confounding with power exceeding the significance level.
\end{theorem}

\begin{proof}[Proof of CD.1]
Proof strategy: construct a measure coupling between World A and World B, prove that $\cI_A$ has the same distribution in both worlds, hence the likelihood ratio is 1 and test power cannot exceed the level.

\textbf{Step 1: Decomposed representation of the Cercis Score.}
Write the Cercis Score as the sum of three components:
\begin{equation}\label{eq:s-decomp-formal}
    S(x) = \underbrace{S_0(x)}_{\text{perfect labels}}
    + \underbrace{\eta \cdot \delta_\varepsilon(x)}_{\text{noise perturbation}}
    + \underbrace{\beta \cdot \delta_z(x)}_{\text{confounding perturbation}},
\end{equation}
where $S_0(x) \defeq Q_0(x) + \eta N_0(x)$ is based on ``ideal'' noise-free, confounding-free labels,
$\delta_\varepsilon(x) \defeq \frac{\partial S}{\partial \varepsilon}(x) \cdot \varepsilon$ is the first-order variation of the Score with respect to noise,
$\delta_z(x) \defeq \frac{\partial S}{\partial z}(x) \cdot z$ is the first-order variation with respect to confounding.
Since $\varepsilon \perp (x,z)$ (Assumption A1), $\delta_\varepsilon$ and $\delta_z$ are conditionally independent given $x$.

\textbf{Step 2: Constructing World B's noise distribution.}
Define World B's noise distribution $P_{\varepsilon'}$ as
\begin{equation}\label{eq:construct-pe}
    P_{\varepsilon'}(A) \defeq \int_{\R \times \R^m}
    \mathbbm{1}_A(\varepsilon + \beta\cdot z)\; dP_\varepsilon(\varepsilon)\, dP_Z(z),\quad
    \forall A \in \cB(\R).
\end{equation}
That is, $P_{\varepsilon'} = \Law(\varepsilon + \beta\cdot z)$ is the marginal distribution of the sum $\varepsilon + \beta\cdot z$.
Since the construction only depends on marginal distributions, $P_{\varepsilon'}$ satisfies:
\begin{itemize}
    \item If $\beta = 0$, then $P_{\varepsilon'} = P_\varepsilon$ (degenerating World B to World A's $H_0$ case).
    \item $\E_{\varepsilon'\sim P_{\varepsilon'}}[\varepsilon'] = \E[\varepsilon] + \beta\cdot\E[z] = 0$, assuming $\E[\varepsilon] = \E[z] = 0$.
    \item $\V[\varepsilon'] = \V[\varepsilon] + \beta^\top \V[z]\beta$ (using $\varepsilon \perp z$, by A1).
\end{itemize}

\textbf{Step 3: Conditional distribution equivalence.}
For any fixed $x \in \cX$, examine the conditional distribution of labels under World A and World B.
Under World A:
\begin{equation}
    P(y \mid x, \text{World A}) = P_\varepsilon\big(y - f^*(x) - \beta\cdot z\big) \cdot P_Z(z).
\end{equation}
Under World B:
\begin{equation}
    P(y \mid x, \text{World B}) = P_{\varepsilon'}\big(y - f^*(x)\big).
\end{equation}

Marginalizing over $z$, the marginal label distribution in World A equals that in World B:
\begin{align}
    P(y \mid x, \text{World A})_{\text{marg}}
    &= \int P_\varepsilon(y - f^*(x) - \beta\cdot z)\, dP_Z(z) \notag\\
    &= P_{\varepsilon'}(y - f^*(x)) \quad\text{(by~\eqref{eq:construct-pe})} \notag\\
    &= P(y \mid x, \text{World B}). \label{eq:marginal-equality}
\end{align}

\textbf{Step 4: Conditional distribution equivalence of Cercis Scores.}
The Cercis Score $S(x) = Q(x) + \eta N(x)$ is a functional of the data distribution. Under finite samples, $S(x)$ depends on the empirical distribution $\hat{P}_{(X,Y)}^{(n)}$. We prove that under $\eta \perp z$, the finite-dimensional distributions of $S(x)$ are identical in both worlds.

Let $S^{(A)}(x)$ and $S^{(B)}(x)$ denote the Cercis Scores in World A and World B respectively.
Using the decomposition~\eqref{eq:s-decomp-formal}:
\begin{align}
    S^{(A)}(x) - S_0(x) &= \eta \cdot \delta_\varepsilon^{(A)}(x) + \beta \cdot \delta_z(x), \\
    S^{(B)}(x) - S_0(x) &= \eta' \cdot \delta_{\varepsilon'}(x),
\end{align}
where $\eta' \indist \eta$. Under the extended independence assumption $\eta \perp (\varepsilon, z)$, the joint distribution of $(\eta, \delta_\varepsilon^{(A)}(x) + \beta\cdot\delta_z(x))$ equals that of $(\eta', \delta_{\varepsilon'}(x))$, yielding identical distributions for $S^{(A)}(x)$ and $S^{(B)}(x)$.

\textbf{Step 5: Likelihood ratio identically 1.}
Let $P_0$ and $P_1$ be the distributions of $\cI_A$ under $H_0$ (World B) and $H_1$ (World A) respectively. By Step 4, $P_0 = P_1$. Hence for any realization $i \in \cI_A$, the likelihood ratio satisfies $\Lambda(i) \defeq dP_1/dP_0(i) = 1$ $P_0$-a.s.

\textbf{Step 6: Power upper bound.}
By the Neyman--Pearson lemma, the most powerful level-$\alpha$ test is based on the likelihood ratio. Since $\Lambda \equiv 1$, any test $\psi$ has power $\E_{H_1}[\psi] = \E_{H_0}[\psi] = \alpha$, establishing the bound with an $O(n^{-1/2})$ finite-sample correction. $\square$
\end{proof}

% ============================================================
\section{Conditional Weak Separability (Theorem CD.2)}
\label{sec:weak}
% ============================================================

\begin{theorem}[Conditional Weak Separability]\label{thm:weak-sep}
When $\I(\eta;z) > 0$ and the Situs encoding preserves the causal skeleton, there exists a curl-based test statistic $\tau$ with non-trivial power for detecting confounding.
\end{theorem}

The key insight is that the Cercis Score gradient $\grad S$ is a conservative field ($\curl(\grad S) = 0$), but kernel-smoothed \textbf{estimators} $\widehat{\grad S}$ need not preserve conservativity --- their non-zero curl components reflect estimation bias and the geometric signature of confounding.

% ============================================================
\section{Causal Relay Bound (Theorem CD.3)}
\label{sec:relay}
% ============================================================

\begin{theorem}[Causal Relay Bound]\label{thm:causal-relay}
Under the front-door criterion and T7 cross-domain fidelity, the Cercis Score shift induced by confounding satisfies
\[
\|\Delta S\|_{\cD_t} \leq L_{\Situs}\cdot\varepsilon + \kappa\cdot\delta + C\cdot\varepsilon\cdot\delta + O(1/\sqrt{n}),
\]
with each term explicitly computable from observable quantities.
\end{theorem}

% ============================================================
\section{Conclusion}
% ============================================================

CD-Theorem establishes that label noise and unobserved confounding are absolutely inseparable when the Sprint parameter is independent of confounders (CD.1), conditionally separable when mutual information exists and the Situs skeleton is preserved (CD.2), and that confounding-induced Cercis bias is bounded by a Situs-distance-dependent expression under the front-door criterion (CD.3).

% ============================================================
\bibliographystyle{plainnat}
\begin{thebibliography}{10}

\bibitem{scx2024cercis}
SCX Working Group.
\newblock ``The \textsf{SCX} Axiomatic Framework: Cercis, Situs, and Tiresias Operators,''
\newblock \emph{Technical Report}, 2024.

\end{thebibliography}

\end{document}
